\documentclass[11pt]{article}

\usepackage{amsmath, amssymb, amsthm}
\usepackage{geometry}
\usepackage{hyperref}
\usepackage{enumitem}
\usepackage{mathtools}
\geometry{margin=1in}

\title{Prime-Indexed Archaeological Dynamics:\\
From Multiplicity Theory to a Testable \texorpdfstring{$\psi$}{psi}-Simulator}
\author{Citizen Gardens \\ \\ The Foundation of Multiplicity}
\date{April 4, 2026}

\begin{document}
\maketitle

\begin{abstract}
This report consolidates a sequence of developments that move from
conceptual Multiplicity Theory to an explicitly testable dynamical
framework for archaeological corpora, with Ancient Egyptian material
as a primary case study. We formalize a prime-indexed state space
\(\psi(t)\) defined on log-transformed prime weights, introduce a
Laplacian-driven stochastic recursion with optional myth operators,
and specify falsifiable observables (entropy, convergence, spectral
structure, and cross-period orbits). A minimal, directly runnable
simulation module is presented in Python, designed to plug into
Egyptian glyph frequency and co-occurrence data. The goal is not
maximal abstraction but a minimum-viable, failure-capable system:
rich enough to encode multiplicity and mythic perturbations, and
simple enough to be decisively supported or refuted by small
corpora.
\end{abstract}
\newpage
\tableofcontents
\newpage
\section{Executive Summary}

\subsection{Context}

Multiplicity Theory reinterprets mathematical and empirical systems
as recursively generated patterns of prime-labeled interactions,
rather than static sets. In the archaeological setting, this means:
\begin{itemize}[nosep]
  \item Each factor in a system (glyph, artifact feature, site property)
        is assigned a unique prime label.
  \item Composite configurations (inscriptions, assemblages, sites) are
        prime factorizations, preserving decomposability and decodability.
  \item Dynamics act on \emph{prime multiplicities} rather than on
        surface tokens, preserving the core ``prime identity'' constraint.
\end{itemize}

Across the thread, several ingredients crystallized:
\begin{enumerate}[nosep]
  \item Prime-based encodings for Egyptian hieroglyphs and archaeological
        factors, with ideas for cross-dialect (ME/LE) recursion,
        entropy tracking, and orbit analysis.
  \item A move from integers themselves to a \emph{log-prime} state
        \(\psi(t)\), compatible with both multiplicative structure and
        numerical stability.
  \item A Laplacian-based recursion driven by a co-occurrence matrix,
        plus localized myth operators that perturb specific prime clusters
        (e.g.\ solar vs.\ chaos glyphs).
  \item A minimal \(\psi\)-simulator that implements these ideas as a
        stochastic linear dynamical system with built-in null models
        and failure criteria.
\end{enumerate}

\subsection{Core Idea}

The culminating framework can be summarized as:

\begin{itemize}[nosep]
  \item \textbf{State:} A gauge-fixed log-prime vector
    \(\psi(t)\in\mathbb{R}^n\) representing the current multiplicity
    of primes (glyphs, features, factors) at epoch \(t\).
  \item \textbf{Geometry:} A co-occurrence matrix \(T\) and its graph
    Laplacian \(L\) encode empirical relational structure.
  \item \textbf{Dynamics:} A discrete-time recursion
    \(\psi(t+1) = \Pi\bigl(\psi(t) - \lambda L\psi(t) + H(t)\psi(t) +
    \sigma \eta(t)\bigr)\), where \(H(t)\) encodes mythic or regime-specific
    perturbations, \(\eta(t)\) is noise, and \(\Pi\) re-centers the state.
  \item \textbf{Observables:} Entropy trajectories, convergence norms,
    spectral properties of \(L\), and cross-period orbit behavior
    constitute the main measurable outputs.
\end{itemize}

Crucially, this system is \emph{falsifiable}: if the empirical geometry
\(T\) does not outperform a shuffled baseline, or if myth operators
have no observable effect, the current multiplicity encoding has
failed to capture real structure.

\subsection{Intended Use}

The immediate application is to Ancient Egyptian corpora, such as
Middle versus Late Egyptian hieroglyphic or hieratic texts, where:

\begin{itemize}[nosep]
  \item Prime labels correspond to glyph types or lexical units.
  \item Weights come from period-specific frequencies.
  \item Co-occurrence matrices come from sliding windows or inscription-level
        co-presence.
  \item Myth operators target curated glyph clusters (solar, chaos, order).
\end{itemize}

The same machinery extends to artifacts, sites, and regional systems
once factors are encoded as primes.

\section{Mathematical Overview}

\subsection{Prime Basis and Raw Weights}

Let \(\mathcal{P} = \{p_1,\dots,p_n\}\) be the prime basis, where each
prime \(p_i\) is associated to a factor \(C_i\). Depending on the
application, factors may be:
\begin{itemize}[nosep]
  \item Hieroglyph types, transliterated signs, or lemma classes.
  \item Artifact features (material, form, function).
  \item Site-level attributes (ritual importance, trade role, etc.).
\end{itemize}

For a given time-slice, corpus, or context \(t\), we collect a raw
weight vector
\[
w(t) = (w_1(t),\dots,w_n(t))^\top,\quad w_i(t)\ge 0,
\]
where \(w_i(t)\) may be counts, frequencies, or other importance scores.

\subsection{Gauge-Fixed Log State}

Directly evolving integers corresponding to prime products is both
numerically fragile and structurally redundant. Instead, we define
a \emph{gauge-fixed log state}:
\[
\psi_i(t)
  = \log\bigl(w_i(t) + \epsilon\bigr)
    - \frac{1}{n}\sum_{k=1}^n \log\bigl(w_k(t) + \epsilon\bigr),
\]
where \(\epsilon>0\) is a small constant ensuring that zero weights
do not cause divergences.

This construction:
\begin{itemize}[nosep]
  \item Preserves prime-indexability: each coordinate corresponds to
        a prime-labeled factor.
  \item Respects multiplicative semantics: multiplicative changes in
        \(w_i\) become additive shifts in \(\psi_i\).
  \item Fixes the gauge: subtracting the mean log-weight removes
        arbitrary uniform offsets, focusing dynamics on \emph{relative}
        prime saliences.
\end{itemize}

We denote \(\psi(t)\in\mathbb{R}^n\) the vector with components
\(\psi_i(t)\).

\subsection{Empirical Geometry: Co-Occurrence and Laplacian}

Empirical relations between factors are captured by a co-occurrence
matrix
\[
T \in \mathbb{R}^{n\times n},\quad T_{ij}\ge 0,
\]
where \(T_{ij}\) measures the strength of association between primes
\(p_i\) and \(p_j\). In textual applications, this may be:
\begin{itemize}[nosep]
  \item Window-based co-occurrence counts of glyphs or tokens.
  \item Inscription-level co-presence frequencies.
  \item Normalized counts (e.g.\ PMI-like or cosine-normalized measures).
\end{itemize}

To obtain a stable spectral geometry, we typically enforce symmetry,
e.g.
\[
T \leftarrow \tfrac{1}{2}\bigl(T + T^\top\bigr),
\]
and optionally apply a normalization such as
\[
\tilde{T} = D^{-1/2} T D^{-1/2},
\]
with \(D = \mathrm{diag}(\sum_j T_{ij})\).

The (combinatorial) graph Laplacian is then
\[
L = D - T,
\]
where \(D_{ii} = \sum_{j} T_{ij}\). The spectral properties of \(L\)
encode the global structure of the co-occurrence graph: connected
components, clusters, and diffusion modes.

\subsection{Dynamics}

\subsubsection{Projection Operator}

Since \(\psi(t)\) is defined only up to an additive constant, it is
natural to enforce zero mean at each step. Define the projection
\[
\Pi(x) = x - \frac{1}{n}\left(\sum_{i=1}^n x_i\right)\mathbf{1},
\]
where \(\mathbf{1}\) is the all-ones vector. This maintains
\(\sum_i \psi_i(t) = 0\) over time.

\subsubsection{Base Recursion}

The base recursion, without mythic perturbations, is
\begin{equation}
\psi(t+1)
= \Pi\Bigl(
    \psi(t) - \lambda L\psi(t) + \sigma \eta(t)
  \Bigr),
\label{eq:base-dynamics}
\end{equation}
where:
\begin{itemize}[nosep]
  \item \(\lambda>0\) is a smoothing parameter.
  \item \(\eta(t)\sim \mathcal{N}(0, I_n)\) is Gaussian noise.
  \item \(\sigma\ge 0\) controls noise magnitude.
\end{itemize}

The term \(-\lambda L\psi(t)\) drives \(\psi\) along the co-occurrence
geometry: roughly, neighbors in the graph pull each other toward
similar values, while high-degree nodes are regularized by the degree
term in \(L\).

In the absence of noise and with \(\lambda\) sufficiently small, we
expect convergence toward modes associated with low Laplacian
eigenvalues (smooth directions over the graph).

\subsubsection{Myth Operators}

Mythic or regime-specific perturbations are modeled by a time-dependent
operator \(H(t)\in\mathbb{R}^{n\times n}\), typically sparse. The full
dynamics is
\begin{equation}
\psi(t+1)
= \Pi\Bigl(
    \psi(t) - \lambda L\psi(t) + H(t)\psi(t) + \sigma\eta(t)
  \Bigr).
\label{eq:full-dynamics}
\end{equation}

The simplest version is diagonal:
\[
(H_{\text{myth}})_{ii} =
\begin{cases}
+\gamma, & i\in S,\\
-\gamma, & i\in C,\\
0, & \text{otherwise},
\end{cases}
\]
where \(S\) and \(C\) are index sets corresponding to, for example,
solar/order glyphs and chaos/serpent glyphs in Egyptian myth, and
\(\gamma>0\) controls perturbation strength.

This turns myth into a localized perturbation in multiplicity space:
\begin{itemize}[nosep]
  \item Primes in \(S\) are systematically amplified.
  \item Primes in \(C\) are systematically suppressed.
\end{itemize}

More elaborate myth operators may include off-diagonal entries
to model conversions (e.g.\ chaos symbols being redirected into
order-linked neighborhoods), but the diagonal form is the minimal
falsifiable starting point.

\subsection{Cross-Period Transfer and Orbits}

A core idea in the Egyptian case is to treat different periods or
dialects (Middle Egyptian, Late Egyptian) as having:
\begin{itemize}[nosep]
  \item The same prime index set \(\mathcal{P}\), but
  \item Different weight estimators and decoding maps.
\end{itemize}

Let:
\begin{itemize}[nosep]
  \item \(\psi^{(\mathrm{ME})}\) be the state corresponding to ME
        weights and decoding.
  \item \(\psi^{(\mathrm{LE})}\) be the state corresponding to LE
        weights and decoding.
\end{itemize}

A cross-period transfer operator \(\mathcal{T}_{\mathrm{ME}\to\mathrm{LE}}\)
can be conceptualized as:
\begin{enumerate}[nosep]
  \item Decode an inscription using the ME mapping to obtain a glyph
        sequence.
  \item Re-encode this sequence under the LE mapping to obtain a new
        weight vector.
  \item Build the corresponding LE-state via the log-map.
\end{enumerate}

Iterating such transfers (ME$\to$LE$\to$ME$\to\cdots$) defines an
orbit in the state space. Orbit length, stability, and attractors
provide concrete measures of semantic drift and structural stability
across periods.

\section{Observables and Failure Modes}

\subsection{Entropy}

From a state \(\psi\), define a probability distribution by a softmax:
\[
\pi_i(\psi) = \frac{\exp(\psi_i)}{\sum_{j=1}^n \exp(\psi_j)}.
\]
The Shannon entropy is:
\[
S(\psi) = -\sum_{i=1}^n \pi_i(\psi)\,\log \pi_i(\psi).
\]

Tracking \(S(t) = S(\psi(t))\) over time reveals:

\begin{itemize}[nosep]
  \item Concentration vs.\ dispersion of prime salience.
  \item The impact of myth operators (e.g.\ entropy spikes or drops
        during mythic phases).
  \item Differences between corpora or periods (ritual vs.\ bureaucratic,
        transitional vs.\ stable).
\end{itemize}

\subsection{Convergence Norm}

Define the convergence norm
\[
c(t) = \|\psi(t+1) - \psi(t)\|_2.
\]
This measures how quickly the system is settling into a fixed point
or low-dimensional attractor under the dynamics.

One can compare convergence rates:
\begin{itemize}[nosep]
  \item Between corpora (e.g.\ ME vs.\ LE).
  \item Between conditions (with and without myth operator).
  \item Against null models (shuffled co-occurrence).
\end{itemize}

\subsection{Spectral Structure}

The eigenvalues and eigenvectors of \(L\) encode the geometry of the
co-occurrence graph. Let
\[
L v_k = \lambda_k v_k,\quad
0 = \lambda_1 \le \lambda_2 \le \cdots \le \lambda_n.
\]

Relevant diagnostics include:
\begin{itemize}[nosep]
  \item \emph{Eigengaps} between small eigenvalues, indicating cluster
        structure or ``semantic basins''.
  \item Alignment of \(\psi(t)\) with leading eigenvectors over time.
  \item Separation between period-specific states (ME vs.\ LE) in the
        subspace spanned by low-\(\lambda_k\) eigenvectors.
\end{itemize}

A strong failure signal is when the spectrum and corresponding
embeddings derived from empirical \(T\) are not meaningfully different
from those derived from a carefully constructed null model.

\subsection{Cross-Period Orbits}

Let \(\mathcal{T}_{\mathrm{ME}\to\mathrm{LE}}\) and
\(\mathcal{T}_{\mathrm{LE}\to\mathrm{ME}}\) denote cross-period
transfer maps on states (abstractly):
\[
\psi^{(\mathrm{ME})}_{k+1}
 = \mathcal{T}_{\mathrm{LE}\to\mathrm{ME}}
   \circ \mathcal{T}_{\mathrm{ME}\to\mathrm{LE}}
   (\psi^{(\mathrm{ME})}_k),
\]
and similarly for LE.

The orbit length and eventual behavior (fixed point, limit cycle,
chaotic-like wandering) define measurable signatures of semantic
stability versus drift. Predictions include:
\begin{itemize}[nosep]
  \item Short cycles in culturally stable regimes.
  \item Longer or more irregular cycles in transitional periods.
\end{itemize}

\subsection{Null Models and Hard Failure Criteria}

A robust null model must preserve some structure (e.g.\ degree
distribution) rather than destroy all geometry. Examples:
\begin{itemize}[nosep]
  \item Row-wise shuffling of \(T\) while keeping row sums fixed.
  \item Degree-preserving edge rewiring on the co-occurrence graph.
\end{itemize}

Hard failure criteria include:

\begin{enumerate}[nosep]
  \item \textbf{No difference to null:} Entropy trajectories,
        convergence norms, and spectral diagnostics for empirical
        \(T\) are statistically indistinguishable from those for
        null \(T_{\text{null}}\).
  \item \textbf{Myth operator inefficacy:} Turning on \(H(t)\) for
        curated clusters (e.g.\ solar vs.\ chaos glyphs) has no
        measurable effect on entropy, convergence, or state
        trajectories.
  \item \textbf{No period separability:} Period-specific states (ME,
        LE, etc.) are not more separable in low-\(\lambda\) eigenvector
        coordinates than appropriate randomized baselines.
\end{enumerate}

If these conditions are met, the current prime-indexed encoding
and dynamics must be reconsidered.

\section{Python \texorpdfstring{$\psi$}{psi}-Simulator Module}

This section presents a minimal, directly usable Python module
implementing the core ideas. It is designed to be:
\begin{itemize}[nosep]
  \item Compatible with arbitrary co-occurrence matrices and weight
        vectors.
  \item Explicit about all parameters.
  \item Equipped with null-model construction and basic diagnostics.
\end{itemize}

\subsection{Code Listing}

\begin{verbatim}
import numpy as np

# -----------------------------
# 1. STATE CONSTRUCTION
# -----------------------------

def build_psi(w, epsilon=1e-8):
    """
    Build gauge-fixed log state psi from raw weights.

    Parameters
    ----------
    w : array-like, shape (n,)
        Raw weight vector (counts, frequencies, etc.).
    epsilon : float
        Small constant to avoid log(0).

    Returns
    -------
    psi : ndarray, shape (n,)
        Centered log-state psi.
    """
    w = np.asarray(w, dtype=float)
    log_w = np.log(w + epsilon)
    psi = log_w - np.mean(log_w)
    return psi


# -----------------------------
# 2. CO-OCCURRENCE → LAPLACIAN
# -----------------------------

def symmetrize(T):
    """
    Symmetrize co-occurrence matrix.
    """
    T = np.asarray(T, dtype=float)
    return 0.5 * (T + T.T)

def normalize_T(T):
    """
    Optional symmetric normalization: D^{-1/2} T D^{-1/2}.
    """
    T = np.asarray(T, dtype=float)
    d = np.sum(T, axis=1)
    D_inv_sqrt = np.diag(1.0 / np.sqrt(d + 1e-8))
    return D_inv_sqrt @ T @ D_inv_sqrt

def build_laplacian(T, normalize=False):
    """
    Build Laplacian L = D - T (optionally on normalized T).

    Parameters
    ----------
    T : ndarray, shape (n, n)
        Co-occurrence matrix.
    normalize : bool
        If True, apply symmetric normalization first.

    Returns
    -------
    L : ndarray, shape (n, n)
        Graph Laplacian.
    """
    T = symmetrize(T)
    if normalize:
        T = normalize_T(T)
    d = np.sum(T, axis=1)
    D = np.diag(d)
    L = D - T
    return L


# -----------------------------
# 3. MYTH OPERATORS
# -----------------------------

def build_myth_operator(n, solar_idx=None, chaos_idx=None, gamma=0.1):
    """
    Build diagonal myth operator.

    Parameters
    ----------
    n : int
        Dimension of state space.
    solar_idx : list of int, optional
        Indices of primes to amplify.
    chaos_idx : list of int, optional
        Indices of primes to suppress.
    gamma : float
        Perturbation magnitude.

    Returns
    -------
    H : ndarray, shape (n, n)
        Myth operator matrix.
    """
    H = np.zeros((n, n), dtype=float)

    if solar_idx is not None:
        for i in solar_idx:
            H[i, i] += gamma

    if chaos_idx is not None:
        for i in chaos_idx:
            H[i, i] -= gamma

    return H


# -----------------------------
# 4. DYNAMICS STEP
# -----------------------------

def project_center(x):
    """
    Project vector to zero-mean subspace.
    """
    x = np.asarray(x, dtype=float)
    return x - np.mean(x)

def step(psi, L, lam=0.1, H=None, sigma=0.0, rng=None):
    """
    One update step:
        psi(t+1) = Pi(psi(t) - lam * L psi(t) + H psi(t) + noise)

    Parameters
    ----------
    psi : ndarray, shape (n,)
        Current state.
    L : ndarray, shape (n, n)
        Laplacian.
    lam : float
        Smoothing parameter lambda.
    H : ndarray, shape (n, n) or None
        Myth operator (optional).
    sigma : float
        Noise scale (Gaussian).
    rng : np.random.Generator or None
        Random number generator.

    Returns
    -------
    psi_next : ndarray, shape (n,)
        Updated state.
    """
    if rng is None:
        rng = np.random.default_rng()

    update = -lam * (L @ psi)

    if H is not None:
        update += H @ psi

    if sigma > 0:
        noise = rng.normal(0.0, sigma, size=psi.shape)
        update += noise

    psi_next = psi + update
    psi_next = project_center(psi_next)
    return psi_next


# -----------------------------
# 5. ENTROPY
# -----------------------------

def compute_entropy(psi):
    """
    Softmax → entropy.
    """
    psi = np.asarray(psi, dtype=float)
    psi_shift = psi - np.max(psi)
    exp_psi = np.exp(psi_shift)
    p = exp_psi / np.sum(exp_psi)
    entropy = -np.sum(p * np.log(p + 1e-12))
    return entropy


# -----------------------------
# 6. SIMULATION
# -----------------------------

def simulate(psi0, L, steps=20, lam=0.1, H=None, sigma=0.0, rng=None):
    """
    Run dynamics for a fixed number of steps.

    Returns
    -------
    out : dict with keys:
        'trajectory' : ndarray, shape (steps+1, n)
        'entropy'    : ndarray, shape (steps+1,)
        'convergence': ndarray, shape (steps,)
    """
    if rng is None:
        rng = np.random.default_rng()

    psi = np.asarray(psi0, dtype=float).copy()
    psi = project_center(psi)

    trajectory = [psi.copy()]
    entropy_series = [compute_entropy(psi)]
    convergence = []

    for _ in range(steps):
        psi_next = step(psi, L, lam=lam, H=H, sigma=sigma, rng=rng)
        c = np.linalg.norm(psi_next - psi)
        convergence.append(c)

        psi = psi_next
        trajectory.append(psi.copy())
        entropy_series.append(compute_entropy(psi))

    return {
        "trajectory": np.array(trajectory),
        "entropy": np.array(entropy_series),
        "convergence": np.array(convergence),
    }


# -----------------------------
# 7. SPECTRAL ANALYSIS
# -----------------------------

def spectrum(L):
    """
    Eigenvalues and eigenvectors of Laplacian.
    """
    L = np.asarray(L, dtype=float)
    vals, vecs = np.linalg.eigh(L)
    return vals, vecs


# -----------------------------
# 8. NULL MODEL (DEGREE-PRESERVING)
# -----------------------------

def shuffled_laplacian(T, rng=None):
    """
    Simple null model: shuffle entries within each row,
    preserving row sums (degree distribution).

    Parameters
    ----------
    T : ndarray, shape (n, n)
        Original co-occurrence matrix.

    Returns
    -------
    L_null : ndarray, shape (n, n)
        Laplacian of shuffled co-occurrence.
    """
    if rng is None:
        rng = np.random.default_rng()
    T = np.asarray(T, dtype=float).copy()
    n = T.shape[0]
    for i in range(n):
        row = T[i, :].copy()
        rng.shuffle(row)
        T[i, :] = row
    L_null = build_laplacian(T, normalize=False)
    return L_null
\end{verbatim}

\subsection{Example Usage Sketch}

\begin{verbatim}
# Example usage (pseudo-code for an Egyptian corpus)

# w_ME, w_LE : glyph counts for Middle / Late Egyptian
# T          : co-occurrence matrix estimated from corpus

psi_ME = build_psi(w_ME)
T_emp = symmetrize(T)
L_emp = build_laplacian(T_emp, normalize=False)

# Null model
L_null = shuffled_laplacian(T_emp)

# Myth operator (e.g. solar vs chaos glyph indices)
solar_idx = [1, 5, 9]
chaos_idx = [2, 6, 10]
H_myth = build_myth_operator(len(w_ME), solar_idx, chaos_idx, gamma=0.2)

# Simulations
res_emp_null = simulate(psi_ME, L_emp, steps=20, lam=0.1, H=None, sigma=0.0)
res_emp_myth = simulate(psi_ME, L_emp, steps=20, lam=0.1, H=H_myth, sigma=0.0)
res_null     = simulate(psi_ME, L_null, steps=20, lam=0.1, H=None, sigma=0.0)

# Compare entropy, convergence, and possibly project trajectory into
# leading eigenvector coordinates for visualization and analysis.
\end{verbatim}

\section{Recommended Next Steps}

\subsection{Egyptian Pilot}

\begin{enumerate}[nosep]
  \item Select a small set (20--50) of inscriptions split across two
        periods (e.g.\ Middle Egyptian vs.\ Late Egyptian).
  \item Build period-specific weight vectors \(w\) and a shared
        co-occurrence matrix \(T\) from these texts.
  \item Run the empirical and null simulations as in the usage sketch.
  \item Design one or two simple yet principled myth operators based
        on curated glyph clusters (solar, chaos, order, funerary).
  \item Evaluate entropy, convergence, and eigenspace separability
        across conditions.
\end{enumerate}

\subsection{Integration with Multiplicity Theory}

Once the empirical pilot reveals non-trivial structure (or fails),
the next layer is:
\begin{itemize}[nosep]
  \item To tie the Laplacian dynamics back to the richer prime-based
        constructions (recursive prime-frequency cipher, hyperprime
        entanglement, etc.).
  \item To treat the empirical \(T\) not only as observed geometry but
        also as a target for generative prime-operator models: the
        ultimate test is whether multiplicity-theoretic generators
        can reproduce the observed spectrum and dynamics of archaeological
        corpora.
\end{itemize}

\section{Conclusion}

The developments summarized here take Multiplicity Theory from a
conceptual proposal about prime-labeled multiplicity spaces to a
concrete, testable dynamical system for archaeological data. The key
features are:
\begin{itemize}[nosep]
  \item A gauge-fixed log-prime state \(\psi(t)\) that respects prime
        decomposability and decodability.
  \item A Laplacian-driven stochastic recursion that encodes empirical
        relational structure.
  \item Myth operators as localized, interpretable perturbations.
  \item A Python module that embodies these ideas with clear null
        models and observational metrics.
\end{itemize}

The strength of this framework lies in its willingness to fail:
if empirical data do not exhibit the predicted signatures, the
encoding and dynamics must be revised. Conversely, if the prime-indexed
system captures robust structure across corpora and periods, it
provides strong evidence that Multiplicity Theory is not just a
narrative overlay but a viable mathematical engine for archaeology,
semiotics, and beyond.

\appendix
\section*{Mathematical Appendix}

\setcounter{section}{0}
\renewcommand{\thesection}{A.\arabic{section}}

\section{State Space, Gauge-Fixing, and Well-Posedness}

\subsection{Prime-indexed state space}

Let \(\mathcal{P} = \{p_1,\dots,p_n\}\) be a fixed index set of primes, each associated with a factor \(C_i\) (glyph, feature, etc.). For a given time-step \(t\), we observe nonnegative weights
\[
w(t) = (w_1(t),\dots,w_n(t))^\top \in [0,\infty)^n.
\]
We assume throughout that
\begin{equation}
\exists\,\epsilon_0>0:\quad w_i(t)+\epsilon \ge \epsilon_0\quad
\text{for all } i,t\text{ and all } \epsilon\in(0,\epsilon_0].
\label{eq:epsilon-lower-bound}
\end{equation}

\subsection{Gauge-fixed log state}

Fix \(\epsilon\in(0,\epsilon_0]\). Define
\[
\ell_i(t) := \log\bigl(w_i(t)+\epsilon\bigr),\qquad
\bar{\ell}(t) := \frac{1}{n}\sum_{k=1}^n \ell_k(t),
\]
and set
\begin{equation}
\psi_i(t)
  := \ell_i(t) - \bar{\ell}(t)
  = \log\bigl(w_i(t)+\epsilon\bigr)
    - \frac{1}{n}\sum_{k=1}^n \log\bigl(w_k(t)+\epsilon\bigr),
\label{eq:psi-def}
\end{equation}
so that
\[
\psi(t) = \bigl(\psi_1(t),\dots,\psi_n(t)\bigr)^\top \in \mathbb{R}^n
\]
and, by construction,
\begin{equation}
\sum_{i=1}^n \psi_i(t) = 0
\quad\text{for all } t.
\label{eq:psi-zero-mean}
\end{equation}
Thus the state space is the codimension-one subspace
\[
\mathcal{H}_0 := \Bigl\{ x\in\mathbb{R}^n : \sum_{i=1}^n x_i = 0 \Bigr\}.
\]

\subsection{Gauge invariance}

\begin{lemma}[Gauge invariance]
\label{lem:gauge-invariance}
Let \(w(t)\) and \(\tilde{w}(t)\) be related by a uniform multiplicative rescaling,
\[
\tilde{w}_i(t) = c(t)\,w_i(t),\qquad c(t) > 0,\ i=1,\dots,n.
\]
Then, for each \(t\), the corresponding gauge-fixed states \(\psi(t)\) and
\(\tilde{\psi}(t)\) defined by \eqref{eq:psi-def} coincide in the limit
\(\epsilon\to 0\) and agree up to \(O(\epsilon)\) for \(\epsilon>0\) small.
\end{lemma}

\begin{proof}
For \(\epsilon>0\), we have
\[
\tilde{\ell}_i(t)
  = \log\bigl(\tilde{w}_i(t)+\epsilon\bigr)
  = \log\bigl(c(t)\,w_i(t) + \epsilon\bigr).
\]
Assuming \(w_i(t)>0\) for the moment, we can rewrite
\[
\tilde{\ell}_i(t)
  = \log\bigl(c(t)\bigr) + \log\Bigl(w_i(t) + \tfrac{\epsilon}{c(t)}\Bigr).
\]
Then
\[
\tilde{\bar{\ell}}(t)
  = \frac{1}{n}\sum_{k=1}^n \tilde{\ell}_k(t)
  = \log\bigl(c(t)\bigr) + \frac{1}{n}\sum_{k=1}^n
    \log\Bigl(w_k(t) + \tfrac{\epsilon}{c(t)}\Bigr).
\]
Hence
\[
\tilde{\psi}_i(t)
  = \tilde{\ell}_i(t) - \tilde{\bar{\ell}}(t)
  = \Bigl[\log\bigl(c(t)\bigr) + \log\Bigl(w_i(t) + \tfrac{\epsilon}{c(t)}\Bigr)\Bigr]
    - \Bigl[\log\bigl(c(t)\bigr) + \frac{1}{n}\sum_{k=1}^n
     \log\Bigl(w_k(t) + \tfrac{\epsilon}{c(t)}\Bigr)\Bigr]
\]
\[
= \log\Bigl(w_i(t) + \tfrac{\epsilon}{c(t)}\Bigr)
  - \frac{1}{n}\sum_{k=1}^n
    \log\Bigl(w_k(t) + \tfrac{\epsilon}{c(t)}\Bigr).
\]
In the limit \(\epsilon\to 0\), we recover
\[
\tilde{\psi}_i(t) \to \log w_i(t) - \frac{1}{n}\sum_{k=1}^n \log w_k(t) = \psi_i(t).
\]
For \(\epsilon>0\) small, a first-order expansion of
\(\log\bigl(w_i(t)+\epsilon/c(t)\bigr)\) in \(\epsilon\) shows that
\(\tilde{\psi}_i(t) = \psi_i(t) + O(\epsilon)\) uniformly over \(i\), provided
\(w_i(t)\) are bounded away from zero. A similar argument applies when some
\(w_i(t)=0\) by invoking \eqref{eq:epsilon-lower-bound}.
\end{proof}

Thus \(\psi(t)\) is invariant (up to a negligible \(\epsilon\)-dependent
correction) under uniform scaling of the weight vector, which matches the
intuitive notion that only \emph{relative} prime intensities matter.

\section{Laplacian Dynamics on the Zero-Mean Subspace}

\subsection{Laplacian construction and basic properties}

Let \(T\in\mathbb{R}^{n\times n}\) be a symmetric nonnegative matrix,
\[
T_{ij} \ge 0,\quad T_{ij} = T_{ji},
\]
and define the degree matrix \(D = \mathrm{diag}(d_1,\dots,d_n)\) with
\[
d_i = \sum_{j=1}^n T_{ij}.
\]
The (combinatorial) Laplacian is
\begin{equation}
L := D - T.
\label{eq:laplacian-def}
\end{equation}

\begin{lemma}[Basic Laplacian properties]
\label{lem:laplacian-basics}
Assume \(T\) is symmetric and nonnegative. Then:
\begin{enumerate}
  \item \(L\) is symmetric and positive semidefinite.
  \item \(L\mathbf{1} = 0\), where \(\mathbf{1} = (1,\dots,1)^\top\).
  \item All eigenvalues \(\lambda_k\) of \(L\) satisfy \(\lambda_k \ge 0\).
\end{enumerate}
\end{lemma}

\begin{proof}
Symmetry follows from symmetry of \(T\). For any \(x\in\mathbb{R}^n\),
\[
x^\top L x = x^\top D x - x^\top T x
  = \sum_{i} d_i x_i^2 - \sum_{i,j} T_{ij} x_i x_j.
\]
By symmetry of \(T\),
\[
\sum_{i,j} T_{ij} x_i x_j
  = \frac{1}{2}\sum_{i,j} T_{ij}(x_i^2 + x_j^2 - (x_i - x_j)^2)
  = \sum_i d_i x_i^2 - \frac{1}{2}\sum_{i,j} T_{ij}(x_i - x_j)^2.
\]
Thus
\[
x^\top L x
  = \sum_i d_i x_i^2 - \Bigl(\sum_i d_i x_i^2 - \tfrac{1}{2}\sum_{i,j} T_{ij}(x_i - x_j)^2\Bigr)
  = \frac{1}{2}\sum_{i,j} T_{ij}(x_i - x_j)^2 \ge 0.
\]
Hence \(L\) is positive semidefinite. The kernel condition is immediate:
\[
(L\mathbf{1})_i = d_i - \sum_j T_{ij} = d_i - d_i = 0.
\]
Positivity of eigenvalues follows from positive semidefiniteness.
\end{proof}

\subsection{Projection and invariant subspace}

Define the orthogonal projection \(\Pi:\mathbb{R}^n\to\mathcal{H}_0\) by
\begin{equation}
\Pi(x) := x - \frac{1}{n}\Bigl(\sum_{i=1}^n x_i\Bigr)\mathbf{1}.
\label{eq:projection-def}
\end{equation}
By construction, \(\Pi\) is linear, symmetric, and idempotent:
\(\Pi^2 = \Pi\), and \(\mathrm{im}(\Pi) = \mathcal{H}_0\).

\begin{lemma}[Invariance of \(\mathcal{H}_0\)]
\label{lem:H0-invariant}
Let \(L\) be a Laplacian as in Lemma~\ref{lem:laplacian-basics}. Then
\(\mathcal{H}_0\) is invariant under \(L\), i.e.\ \(x\in\mathcal{H}_0\)
implies \(Lx\in\mathcal{H}_0\). Moreover, \(L\) and \(\Pi\) commute:
\[
L\Pi = \Pi L.
\]
\end{lemma}

\begin{proof}
For \(x\in\mathcal{H}_0\), we have \(\sum_i x_i = 0\). Since \(L\mathbf{1}=0\),
\[
\sum_i (Lx)_i
  = \mathbf{1}^\top L x
  = (L^\top \mathbf{1})^\top x
  = (L\mathbf{1})^\top x = 0.
\]
Thus \(Lx\in\mathcal{H}_0\). For commuting, note that for any \(x\),
\[
L\Pi x
  = L\Bigl(x - \frac{1}{n}(\mathbf{1}^\top x)\mathbf{1}\Bigr)
  = Lx - \frac{1}{n}(\mathbf{1}^\top x)L\mathbf{1}
  = Lx,
\]
while
\[
\Pi L x
  = Lx - \frac{1}{n}(\mathbf{1}^\top Lx)\mathbf{1}
  = Lx - \frac{1}{n}( (L^\top \mathbf{1})^\top x)\mathbf{1}
  = Lx.
\]
So \(L\Pi x = \Pi L x\).
\end{proof}

\section{Deterministic Dynamics: Operator Norms and Stability}

\subsection{Linear part without myth operator}

Consider the deterministic dynamics without mythic perturbation or noise:
\begin{equation}
\psi(t+1) = \Pi\bigl(\psi(t) - \lambda L\psi(t)\bigr),
\label{eq:deterministic-noH}
\end{equation}
which we may write as
\[
\psi(t+1) = A \psi(t),
\]
with
\[
A := \Pi(I - \lambda L).
\]
By Lemma~\ref{lem:H0-invariant}, if \(\psi(0)\in\mathcal{H}_0\), then
\(\psi(t)\in\mathcal{H}_0\) for all \(t\), and we may restrict attention
to \(\mathcal{H}_0\).

Let the eigenvalues of \(L\) be
\[
0 = \lambda_1 \le \lambda_2 \le \cdots \le \lambda_n.
\]
On \(\mathcal{H}_0\), the \(\lambda_1=0\) eigenvalue corresponds to the
constant vector \(\mathbf{1}\), which is projected out. Thus the relevant
spectrum on \(\mathcal{H}_0\) is \(\{\lambda_2,\dots,\lambda_n\}\).

\begin{lemma}[Spectrum of \(A\) on \(\mathcal{H}_0\)]
\label{lem:A-spectrum}
On \(\mathcal{H}_0\), the operator \(A\) is diagonalizable with eigenvalues
\[
\mu_k = 1 - \lambda \lambda_k,\quad k=2,\dots,n.
\]
\end{lemma}

\begin{proof}
Since \(L\) is symmetric, its eigenvectors \(v_k\) form an orthonormal basis.
We have \(L v_k = \lambda_k v_k\). The eigenvector corresponding to
\(\lambda_1=0\) is proportional to \(\mathbf{1}\), which lies outside
\(\mathcal{H}_0\). All other eigenvectors \(v_k\) with \(k\ge 2\) satisfy
\(\sum_i (v_k)_i = 0\), hence lie in \(\mathcal{H}_0\). On such a
\(v_k\),
\[
\Pi v_k = v_k, \qquad
A v_k = \Pi(I - \lambda L) v_k
      = \Pi\bigl((1 - \lambda \lambda_k) v_k\bigr)
      = (1 - \lambda \lambda_k) v_k.
\]
Thus \(\mu_k = 1 - \lambda \lambda_k\) for \(k=2,\dots,n\).
\end{proof}

\subsection{Operator norm bound and contractivity}

Let \(\|\cdot\|\) denote the Euclidean norm on \(\mathcal{H}_0\) and its
induced operator norm. From Lemma~\ref{lem:A-spectrum}, the spectral radius
of \(A\) on \(\mathcal{H}_0\) is
\[
\rho(A) = \max_{k=2,\dots,n} |1 - \lambda \lambda_k|.
\]

\begin{proposition}[Contractivity condition]
\label{prop:contractivity}
Assume \(\lambda>0\) is chosen such that
\[
0 < \lambda < \frac{2}{\lambda_n}.
\]
Then the operator norm of \(A\) on \(\mathcal{H}_0\) satisfies
\[
\|A\| \le \max\bigl\{|1 - \lambda\lambda_2|,\ |1 - \lambda\lambda_n|\bigr\} < 1,
\]
and \eqref{eq:deterministic-noH} is a strict contraction on
\(\mathcal{H}_0\).
\end{proposition}

\begin{proof}
Because \(L\) is symmetric and positive semidefinite, its eigenvalues
are real and nonnegative. On \(\mathcal{H}_0\) we have \(\lambda_k>0\)
for \(k\ge 2\). For a fixed \(\lambda>0\), the function
\(f(\lambda_k) = |1 - \lambda\lambda_k|\) on \([\lambda_2,\lambda_n]\)
attains its maximum at one of the endpoints unless it crosses zero.
If \(0<\lambda< 2/\lambda_n\), then
\[
1 - \lambda\lambda_n > -1,
\quad
1 - \lambda\lambda_k > -1\ \text{for all }k,
\]
so each eigenvalue \(\mu_k = 1 - \lambda\lambda_k\) lies in \((-1,1)\),
and hence \(\|A\| = \rho(A) = \max_k |\mu_k| < 1\). Therefore \(A\) is
a strict contraction on \(\mathcal{H}_0\).
\end{proof}

\begin{corollary}[Exponential convergence]
\label{cor:exp-convergence}
Under the assumptions of Proposition~\ref{prop:contractivity}, the
deterministic dynamics \eqref{eq:deterministic-noH} satisfies
\[
\|\psi(t)\| \le \|A\|^t \,\|\psi(0)\|
\]
for all \(t\ge 0\), and hence converges exponentially to \(0\) in
\(\mathcal{H}_0\). In particular, the system converges to a constant
(log-weight) configuration modulo the zero-mean gauge.
\end{corollary}

\section{Inclusion of Myth Operator}

\subsection{Diagonal myth operator and stability}

Consider now the linear dynamics with a time-independent myth operator
and no noise:
\begin{equation}
\psi(t+1) = \Pi\bigl(\psi(t) - \lambda L\psi(t) + H\psi(t)\bigr)
         = B \psi(t),
\label{eq:deterministic-withH}
\end{equation}
with
\[
B := \Pi(I - \lambda L + H).
\]
We assume throughout that \(H\) is symmetric and that \(\mathbf{1}\) is
an eigenvector of \(H\) (e.g.\ for diagonal \(H\) with \(\sum_i H_{ii}=0\),
this is automatically satisfied after projection), so that \(\mathcal{H}_0\)
remains invariant under \(I - \lambda L + H\) and hence under \(B\).

The diagonal myth operator used in the main text has the form
\[
H_{ii} =
\begin{cases}
+\gamma, & i\in S,\\
-\gamma, & i\in C,\\
0, & \text{otherwise},
\end{cases}
\]
with \(S,C\subset\{1,\dots,n\}\) disjoint and \(\gamma>0\). This is a bounded
self-adjoint operator on \(\mathbb{R}^n\) with operator norm
\(\|H\|\le \gamma\).

\begin{proposition}[Perturbative stability under small myth operator]
\label{prop:myth-stability}
Assume the conditions of Proposition~\ref{prop:contractivity} and that
\(\|H\|\) is sufficiently small so that
\[
\|H\| < 1 - \|A\|.
\]
Then \(B\) is a strict contraction on \(\mathcal{H}_0\), and the dynamics
\eqref{eq:deterministic-withH} is exponentially stable.
\end{proposition}

\begin{proof}
On \(\mathcal{H}_0\), we have
\[
B = A + \Pi H.
\]
Note that \(\|\Pi H\| \le \|\Pi\|\|H\| = \|H\|\). Then
\[
\|B\| \le \|A\| + \|\Pi H\| \le \|A\| + \|H\|.
\]
By assumption \(\|H\| < 1 - \|A\|\), so \(\|B\| < 1\). Hence \(B\)
is a strict contraction on \(\mathcal{H}_0\). The same contraction
argument as in Corollary~\ref{cor:exp-convergence} yields exponential
convergence to a unique fixed point in \(\mathcal{H}_0\).
\end{proof}

\subsection{Fixed point equation}

The fixed point \(\psi^\ast\in\mathcal{H}_0\) satisfies
\[
\psi^\ast = \Pi\bigl(\psi^\ast - \lambda L\psi^\ast + H\psi^\ast\bigr),
\]
or equivalently, since \(\psi^\ast\in\mathcal{H}_0\) and \(L\) commutes
with \(\Pi\),
\[
\psi^\ast = \psi^\ast - \lambda L\psi^\ast + H\psi^\ast - \frac{1}{n}
\Bigl(\sum_{i=1}^n (H\psi^\ast)_i\Bigr)\mathbf{1}.
\]
Projection annihilates the constant component, so within \(\mathcal{H}_0\)
we may write succinctly
\[
(\lambda L - H)\psi^\ast = 0,\quad \psi^\ast \in \mathcal{H}_0.
\]
If \(\lambda L - H\) is invertible on \(\mathcal{H}_0\), the only solution
is \(\psi^\ast = 0\). Nontrivial fixed points thus correspond to eigenvectors
of \(L^{-1}H\) in \(\mathcal{H}_0\). In the small-\(H\) regime of
Proposition~\ref{prop:myth-stability}, any such nontrivial solutions are
ruled out, and the fixed point is unique.

\section{Stochastic Dynamics and Moment Bounds}

\subsection{Stochastic recursion}

The full stochastic recursion considered in the main text is
\begin{equation}
\psi(t+1)
  = \Pi\bigl(\psi(t) - \lambda L\psi(t) + H\psi(t) + \sigma\eta(t)\bigr),
\label{eq:stochastic-dynamics}
\end{equation}
where \(\{\eta(t)\}_{t\ge 0}\) are i.i.d.\ standard Gaussian vectors in
\(\mathbb{R}^n\) and \(\sigma\ge 0\) is a noise amplitude. Let
\[
B := \Pi(I - \lambda L + H)
\]
as before. Then we may write
\[
\psi(t+1) = B\psi(t) + \sigma \xi(t),
\]
where \(\xi(t) := \Pi\eta(t)\) are i.i.d.\ centered Gaussians supported
on \(\mathcal{H}_0\) with covariance \(\Pi\).

\subsection{Second moment bound under contractivity}

Assume that the contractivity condition of Proposition~\ref{prop:myth-stability}
holds, i.e.\ \(\|B\|<1\) on \(\mathcal{H}_0\). Then the recursion is a
linear Gaussian system on \(\mathcal{H}_0\).

\begin{proposition}[Uniform second moment bound]
\label{prop:second-moment}
Under the contractivity assumption \(\|B\|<1\), for any initial state
\(\psi(0)\in\mathcal{H}_0\),
\[
\sup_{t\ge 0} \mathbb{E}\,\|\psi(t)\|^2 < \infty.
\]
More precisely,
\[
\mathbb{E}\,\|\psi(t)\|^2
  \le \|B\|^{2t} \|\psi(0)\|^2
      + \frac{\sigma^2 \mathrm{tr}(\Pi)}{1 - \|B\|^2}.
\]
\end{proposition}

\begin{proof}
We have
\[
\psi(t) = B^t \psi(0) + \sigma \sum_{k=0}^{t-1} B^k \xi(t-1-k),
\]
and the \(\xi(\cdot)\) are independent with mean zero. Then
\[
\mathbb{E}\,\|\psi(t)\|^2
  = \|B^t \psi(0)\|^2
    + \sigma^2 \sum_{k=0}^{t-1} \mathbb{E}\,\|B^k \xi(0)\|^2.
\]
Since \(\|\xi(0)\|^2\) has expectation \(\mathrm{tr}(\Pi) = n-1\) and
\(\|B^k\|\le \|B\|^k\), we get
\[
\mathbb{E}\,\|B^k \xi(0)\|^2
 \le \|B^k\|^2 \mathbb{E}\,\|\xi(0)\|^2
 \le \|B\|^{2k} (n-1).
\]
Thus
\[
\mathbb{E}\,\|\psi(t)\|^2
  \le \|B\|^{2t} \|\psi(0)\|^2
     + \sigma^2 (n-1)\sum_{k=0}^{t-1} \|B\|^{2k}
  \le \|B\|^{2t} \|\psi(0)\|^2
     + \frac{\sigma^2 (n-1)}{1 - \|B\|^2}.
\]
Since \(\|B\|<1\), this is uniformly bounded in \(t\).
\end{proof}

The bound shows that in the contractive regime, the stochastic dynamics
admits a unique stationary distribution on \(\mathcal{H}_0\) with finite
second moment, and trajectories forget their initial conditions
exponentially fast.

\section{Softmax Entropy and Lipschitz Estimates}

\subsection{Definition and basic properties}

Given \(\psi\in\mathcal{H}_0\), define
\[
\pi_i(\psi) = \frac{\exp(\psi_i)}{\sum_{j=1}^n \exp(\psi_j)},\qquad
S(\psi) = -\sum_{i=1}^n \pi_i(\psi) \log \pi_i(\psi).
\]
The map \(\psi\mapsto \pi(\psi)\) is the softmax, and \(S(\psi)\) is
the Shannon entropy of the resulting distribution.

\begin{lemma}[Boundedness and invariance under translations]
\label{lem:entropy-bounds}
For all \(\psi\in\mathbb{R}^n\),
\begin{enumerate}
  \item \(0 \le S(\psi) \le \log n\).
  \item For any constant \(c\in\mathbb{R}\),
        \(S(\psi + c\mathbf{1}) = S(\psi)\).
\end{enumerate}
\end{lemma}

\begin{proof}
The bounds are standard: entropy is minimized (0) when one component is 1
and the rest 0, and maximized (\(\log n\)) when all components are equal.
For translation invariance, note that
\[
\pi_i(\psi + c\mathbf{1}) = \frac{e^{\psi_i + c}}{\sum_j e^{\psi_j + c}}
                          = \frac{e^{\psi_i}}{\sum_j e^{\psi_j}}
                          = \pi_i(\psi),
\]
hence \(S(\psi + c\mathbf{1}) = S(\psi)\).
\end{proof}

Since \(\psi(t)\in\mathcal{H}_0\) are gauge-fixed up to the zero-mean
constraint, the entropy observable used in the main text is well-defined
modulo the natural translation invariance.

\subsection{Local Lipschitz continuity on bounded sets}

While \(S(\psi)\) is not globally Lipschitz on \(\mathbb{R}^n\) with
respect to the Euclidean norm, it is locally Lipschitz on bounded sets.

\begin{proposition}[Local Lipschitz continuity]
\label{prop:entropy-lipschitz}
For any radius \(R>0\), there exists \(L_R>0\) such that for all
\(\psi,\phi\in\mathcal{H}_0\) with \(\|\psi\|\le R\) and \(\|\phi\|\le R\),
\[
|S(\psi) - S(\phi)| \le L_R \|\psi - \phi\|.
\]
\end{proposition}

\begin{proof}
The softmax map \(\psi\mapsto \pi(\psi)\) is smooth on \(\mathbb{R}^n\), and
its Jacobian is bounded on any compact set. Similarly, the function
\(\pi\mapsto S(\pi)\) is smooth on the interior of the simplex. Since the
image of a bounded subset of \(\mathbb{R}^n\) under softmax is contained in
a compact subset of the probability simplex with all coordinates bounded
away from 0, the composition is smooth with bounded derivative there.
Hence, the gradient \(\nabla S(\psi)\) is bounded on
\(\{\psi \in \mathcal{H}_0 : \|\psi\|\le R\}\). Let
\[
M_R := \sup_{\|\psi\|\le R} \|\nabla S(\psi)\| < \infty.
\]
By the mean value theorem,
\[
|S(\psi) - S(\phi)| \le M_R \|\psi - \phi\|
\]
for all \(\psi,\phi\) in the ball. Take \(L_R := M_R\).
\end{proof}

In particular, under the contractive dynamics of
Proposition~\ref{prop:myth-stability} and the moment bound
in Proposition~\ref{prop:second-moment}, trajectories almost surely
enter and remain in a bounded region of \(\mathcal{H}_0\), so entropy
differences can be controlled by differences in the state norm.

\section{Summary of Operator Norm Conditions}

For clarity, we summarize the main operator norm conditions derived
in this appendix:

\begin{itemize}
  \item \textbf{Base Laplacian dynamics:} On \(\mathcal{H}_0\), the
        linear operator \(A = \Pi(I - \lambda L)\) has eigenvalues
        \(1 - \lambda\lambda_k\) for \(k\ge 2\). A sufficient condition
        for strict contractivity is
        \[
        0 < \lambda < \frac{2}{\lambda_n},
        \]
        in which case \(\|A\|<1\).
  \item \textbf{Myth perturbation:} If \(H\) is symmetric and
        \(\|H\| < 1 - \|A\|\), then
        \[
        B = \Pi(I - \lambda L + H)
        \]
        remains a strict contraction on \(\mathcal{H}_0\), and the
        deterministic dynamics with myth operator is exponentially stable.
  \item \textbf{Stochastic stability:} Under the same contractivity
        condition, the stochastic recursion with Gaussian noise has
        bounded second moments uniformly in time, and admits a unique
        stationary distribution on \(\mathcal{H}_0\) with finite variance.
\end{itemize}

These bounds justify the use of the Laplacian-based \(\psi\)-dynamics
as a well-posed and stable backbone for the empirical experiments
outlined in the main body of the article.

\nocite{*}
\bibliographystyle{plainnat}
\bibliography{references}
\end{document}